% !TEX root = ../main.tex
\section{보관해 둔 확장 기준선}
\label{ch:appendix-archive}

같은 매칭 코호트(모든 점수가 똑같이 점수를 매기는 지갑 묶음)와 같은 지표 가족(비슷한 지표 묶음)으로 다른 페이지랭크(PageRank) 변형들도 점수를 매겨 보았어요. 이것들은 \ref{ch:results}장에서 평가한 점수에는 들어 있지 않아요.

\dissertationSubheading[sec:seven-method-archive]{방법 일곱 개의 일치도 행렬}

평가 보관함에는 방법 일곱 개로 만든 행렬(방법과 지표를 엮은 결과 표)이 남아 있어요(LP-PR, CW-AWP, LF-PR, 그리고 RiskProp 방식의 변형들). 하지만 여기에는 그리지 않았어요. 이 행렬의 행들은 C-PR(결합 페이지랭크)에서 금액을 그대로 쓴 두 층을 하나씩 따로 걸은 것을 허락 기준과 송금 기준으로 삼아 계산했어요. 이 방법들 가운데 허락 가족에서는 허락 걷기가 앞서고, 송금 가족에서는 송금 걷기가 그와 아주 가까운 변형들과 함께 앞서요.

보관해 둔 행렬에서는 다음과 같은 모습이 보여요.

\begin{itemize}
\item 송금과 비슷한 화살표를 쓰는 방법들은 송금 가족과 시빌(가짜 주소를 많이 만들어 속이는 공격) 가족에서 송금 걷기 둘레에 모여요.
\item 화살표가 아주 적은 방법은 순위가 불안정하거나 퇴화할(거의 모두 같은 값이 되어 쓸모없어질) 수 있어요. 그리고 실행 시간이 짧은 것도 그저 $m$이 아주 작기 때문일 수 있어요.
\item 숫자가 송금 걷기와 똑같이 나오는 RiskProp 방식의 행들은 같은 송금 구조를 함께 쓴다는 뜻이에요. 따로 무언가를 새로 찾아낸 것이 아니에요.
\end{itemize}

\dissertationSubheading[sec:six-aave-archive]{방법 여섯 개의 Aave 중심 행렬}

Aave(돈을 빌려주고 빌리는 디파이 서비스)에 맞춘 미리 정한 설정 묶음(프리셋)도 보관해 두었어요. 연구 기간 동안 아비트럼(Arbitrum)의 빌려주기 이벤트 그래프 가운데 몇 개는 화살표가 듬성듬성해요. 그리고 \texttt{BorrowAllowanceDelegated} 이벤트를 꺼내지 않았을 때는 위임에 바탕을 둔 행이 퇴화할 수 있어요(부록~\ref{ch:appendix-eval}).

이 프리셋은 보관함에 남아 있고, 여기에는 그리지 않았어요.

이 행렬들은 신용 위임(빌릴 수 있는 권한을 남에게 맡기는 일) 화살표와, 여러 서비스에 걸친 대출 라벨(나중에 채점에 쓰는 정답 값)에 대한 앞으로의 연구를 이끌어 내요. 하지만 주장(C) C1--C5는 바꾸지 않아요.

\dissertationSubheading[sec:why-archived]{이 변형들을 보관만 해 둔 까닭}

평가한 점수가 엔도스랭크(EndorseRank)와 AWP, 그리고 같은 두 층으로 만든 C-PR과 S-PR인 까닭은 세 가지예요.

\begin{enumerate}
\item 허락에 바탕을 둔 보증(믿고 맡김)이 새로운 구성개념(점수가 재려고 하는 생각)이에요. 이것을 이미 자리 잡은 송금 방법\cend{3} 옆에 놓고 평가하면, 연구 질문만 따로 떼어 볼 수 있어요.
\item AWP와 송금 화살표를 함께 쓰는 변형들은 표에 행을 더 보탤 뿐, 따로 선 구성개념을 보태지는 않아요.
\item 연구 계획서에 있던 GF-PR 참고 점수는 결과(거래 성과)에서 바로 만든 점수예요. 이 점수는 거래 성공 대리 지표(대신 재는 값)와 같은 GMX 이벤트에서 나와요. 그래서 둘 사이의 일치는 검증이 아니라 구성개념 겹침으로 쳐요(\ref{sub:gf-pr}절). GF-PR을 넣은 방법 세 개 행렬은 방법 일곱 개 프리셋, Aave 여섯 개 프리셋과 함께 보관해 두었어요.
\end{enumerate}

\dissertationSubheading[sec:archive-repro]{보관 표 다시 만들기}

방법 일곱 개 프리셋과 Aave 여섯 개 프리셋으로 평가 처리 과정의 표 내보내기를 돌리면 보관 표가 다시 만들어져요. 표 조각과 함께 CSV 사본도 저장돼요.

\dissertationSubheading[sec:archive-method-notes]{보관한 방법 이름표에 대한 메모}

보관 표와 CSV로 내보낸 파일은 아래 이름표(식별자)를 써요.

\begin{description}
\item[LP-PR] 금액을 그대로 쓴 송금 층에 AWP의 시간 감쇠를 쓴 것이에요. 화살표 무게마다 1에서 보낸 쪽의 GMX 청산(손실이 너무 커져 거래가 강제로 끝나는 일) 비율을 뺀 값과, 1에서 받는 쪽의 청산 비율을 뺀 값을 곱해요. 송금 가족들에서는 송금 걷기를 바짝 따라가요. 무게에 청산 대리 지표와 같은 청산 결과를 쓰므로, GF-PR과 같은 순환 문제(맞혀야 할 답을 재료로 쓰는 문제)를 함께 가지고 있어요.
\item[CW-AWP] 금액을 그대로 쓴 송금 층을, 미리 정한 담보 토큰 목록의 송금으로만 줄인 것이에요. 송금 $\tau$는 전체 층의 값보다 조금 낮아요.
\item[LF-PR] Aave 이벤트로 만든 대출 흐름 그래프예요. 빌리기는 풀(돈을 모아 둔 곳)에서 사용자로 가는 화살표이고, 갚기와 청산은 사용자에서 풀로 가는 화살표예요. 모두 AWP의 시간 감쇠를 써요. 이 기간에 아비트럼에서 Aave를 쓴 코호트 지갑은 적었어요. 그래서 이 그래프는 화살표가 아주 적고 실행 시간이 거의 0이며, 그 순위는 송금 걷기들과 비교할 수 없어요.
\item[RiskProp] 금액을 그대로 쓴 송금 층에서, 화살표 무게마다 1에서 보낸 쪽의 GMX 청산 비율을 뺀 값을 곱한 것이에요. 위험을 퍼뜨린다는 생각은 Lin 등\cend{12}에게서 빌려 왔지만, 그들의 방법은 아니에요. 이 코호트에서는 여러 가족에서 숫자가 송금 걷기와 거울처럼 똑같이 나와요. 그러니 청산으로 무게를 준 것은 거의 아무것도 바꾸지 않았어요.
\item[LiqCall-PR / Borrow-PR / Repay-PR] 쓸 수 있는 경우에, Aave의 청산 호출, 빌리기, 갚기로 만든 지갑과 지갑 사이의 그래프예요. 연구 기간의 아비트럼에서 이 그래프들은 화살표가 수십 개에서 수백 개뿐이에요. 그래서 일치도는 약하고, 실행 시간은 무시해도 될 만큼 짧아요.
\item[Delegation-PR] 신용 위임 허락을 쓰려고 만든 방법이에요. 꺼낸 위임 화살표가 하나도 없어서, 점수가 퇴화했고 $\tau$도 해석할 수 없어요.
\end{description}

보관한 행 가운데 동료 심사(다른 연구자들의 검토)를 거친 기준선(비교 기준)은 하나도 없어요.

\dissertationSubheading[sec:archive-decision-log]{평가할 점수를 고른 기준}

\ref{ch:results}--\ref{ch:discussion}장은 엔도스랭크, AWP, 그리고 이 둘의 두 층을 결합한 걷기에 집중해요. 다른 변형들은 다음 기준에 미치지 못해요.

\begin{enumerate}
\item 구성개념의 독립성: 기준선은 송금에 무게만 다시 매기는 것이 아니라, 화살표에 다른 뜻을 담아야 해요.
\item 자료의 충분함: 매칭 코호트에서 페이지랭크가 안정적일 만큼 화살표가 충분히 있어야 해요.
\item 범위: 이 연구의 기여(이 연구가 새로 보탠 것)는 허락에 바탕을 둔 보증이에요. 가능한 모든 페이지랭크 변형을 두루 살펴보려는 것이 아니에요.
\end{enumerate}

이 기준에 따르면, 엔도스랭크(허락)와 AWP(송금), 그리고 두 층을 결합한 걷기들이 꼭 필요한 만큼만 갖춘 충분한 묶음이에요. 방법 일곱 개 행렬과 Aave 여섯 개 행렬은 보관된 채로 남고, 주장~C1--C5에는 들어가지 않아요.

